\exnumber{01}
\extitle{UX II - Ex 2 - Mener les entretiens}
\exfiles{ex01/PDF, nomme Name\_FirstName\_DeliverableName\_Date\_VersionNumber}
\exauthorize{Those named in the exercise}
\makeheaderfiles

\textbf{Instructions :}

Run 6 interviews. Go in pairs, one being the interviewer and the other the note taker.
Go outside! Meet random people in cafés, at the bus, on city squares, in parks, on benches, introduce yourself and ask for 15 minutes of their time. Run your interview guide and write down anything that could be helpful, like reactions, hesitations, etc. You can use a dictaphone to make sure you're not missing anything. But in this case, let the interviewee know and ask for their permission!
You can offer them a coffee to thank them, sometimes a smile suffices!
At least 3 of your interviews must be with strangers. You'd be surprised how people are happy to help and looooove talking about themselves.

Then, write down the script on a written document, and give elements of context regarding each interview (place, time, name, etc.)

\textbf{Livrables :}

\begin{itemize}
\item A PDF document with the scripts from all the 6 interviews.
\item Déposez le PDF sur Github, nommé ainsi : « Name\_FirstName\_DeliverableName\_Date\_VersionNumber ». Exemple : « Doe\_John\_MidFiPrototype\_05242024\_V1 ».
\item Ajoutez un titre dans le document.
\item Soignez l'aspect global du document. Il doit être clair, simple et lisible. Vous pouvez regarder canva.com, sans surcharger le dessin de détails.
\end{itemize}

